%===============================================================================
\section{Results IV: Robustness and the Choice of $K$}\label{sec:res-robust}
%===============================================================================

\paragraph{Attacker seeds and data splits}
Ground truth is itself an estimate: attackers are trained with random initializations and
early-stopping splits, and $\Vy$ is measured on one random test split. We repeated the complete
three-attacker retraining with two further seeds (RTS-GMLC, three targets; PJM load) and on two
further random 70/30 splits (RTS-GMLC). Across these reruns the mean absolute change in $\Vy$ is
0.001--0.008, the Spearman correlation of $M^{(2)}$ across fields with the main run is 0.983--0.999,
the set of $\tau$-critical fields is identical at $\tau=0.5$ (Jaccard 1.0 in every case) and agrees
with Jaccard 0.93--1.00 at $\tau=0.7$, and the number of combination-critical fields and the minimum
withholding set change by at most one or two fields (Fig.~\ref{fig:robust}c). The conclusions of
Sections~\ref{sec:res-risk}--\ref{sec:res-def} are therefore not artefacts of a particular seed or
split.

\paragraph{Threshold}
Fig.~\ref{fig:robust}(d) sweeps $\tau$ from 0.3 to 0.9 over all ten targets. At every threshold the
share of fields that are safe alone but critical with at most two background fields exceeds the share
of fields that are unsafe alone (e.g.\ 43\% versus 10\% at $\tau=0.7$), and withholding the fields
flagged by single-field grading leaves between 38\% (at $\tau=0.9$) and 96\% (at $\tau=0.5$) of the
minimal unsafe sets intact. The gap between the two grading philosophies is thus not tied to the
choice of $\tau$.

\paragraph{Choice of $K$}
Exhaustive retraining of all sets of four fields on every data set (for NEM an additional 27,405 sets)
gives the exact $M^{(3)}$. Measured by the \emph{marginal} risk, $K=2$ is close to saturation:
$M^{(2)}/M^{(3)}$, the ratio of mean field risk, lies between 0.958 and 1.000 for nine targets and is
0.813 for NEM TUMUT3, a hydro unit whose dispatch depends on inter-regional exchange
(Fig.~\ref{fig:robust}a). The first background field brings most of the escalation ($M^{(1)}/M^{(2)}$
between 0.73 and 0.97).

Threshold crossings, however, do not saturate as quickly, which is the caveat stated in
Section~\ref{sec:mk}. At $\tau=0.7$ several targets have fields that become critical only with a third
background field: PJM net interchange (3 critical fields at $K=2$, 21 at $K=3$), NEM PPCCGT (10 and
24) and PJM load (14 and 22) (Fig.~\ref{fig:robust}b). At $\tau=0.5$ all fields are already critical at
$K=2$ except for TUMUT3 (22 versus 30). Minimal unsafe sets of size four therefore exist that a
hitting set computed at $K=2$ does not break. Breaking all minimal unsafe sets up to size four
requires on average 11.3 withheld fields instead of 9.4 at $\tau=0.5$, and 6.3 instead of 5.1 at
$\tau=0.7$. Our recommendation follows: $K=2$ suffices to \emph{grade} fields by their marginal risk,
while \emph{disposal} decisions should be computed at $K=3$, at the price of one to two additional
withheld fields. Evaluating all sets of size four is exactly the workload for which the amortized
scan is needed: for CAISO it comprises 52,360 sets per target, against 7,175 for $K=2$.

\begin{figure}[!tbp]
\centering
\includegraphics[width=\textwidth]{fig9_robustness.pdf}
\caption{Robustness. (a) Saturation of mean field risk in the background budget. (b) Share of
$\tau$-critical fields at $\tau=0.7$ with $K=2$ (filled) and $K=3$ (open). (c) Agreement of reruns with
different attacker seeds or data splits with the main run. (d) Threshold sweep over ten targets:
fields that are safe alone but critical (solid), fields unsafe alone (dashed), and minimal unsafe sets
left after single-field grading (dotted).}
\label{fig:robust}
\end{figure}
